\documentclass[11pt,a4paper]{article}

% ------------------------------------------------------------------
% Paket
% ------------------------------------------------------------------
\usepackage[utf8]{inputenc}
\usepackage[T1]{fontenc}
\usepackage{babel}
\babelprovide[import,main]{indonesian}
\usepackage{lmodern}
\usepackage[margin=2.5cm]{geometry}
\usepackage{amsmath,amssymb,amsthm}
\usepackage{physics}
\usepackage{graphicx}
\usepackage{booktabs}
\usepackage{xcolor}
\usepackage{listings}
\usepackage{caption}
\usepackage{enumitem}
\usepackage[hidelinks]{hyperref}

\newtheorem{teorema}{Teorema}
\newtheorem{definisi}{Definisi}
\newtheorem{lema}{Lema}

% Penanda bagian yang masih harus diisi
\newcommand{\todo}[1]{\textcolor{red}{\textbf{[TODO:} #1\textbf{]}}}

\lstset{
  basicstyle=\ttfamily\small,
  keywordstyle=\color{blue!70!black}\bfseries,
  commentstyle=\color{green!40!black}\itshape,
  stringstyle=\color{orange!70!black},
  numbers=left, numberstyle=\tiny\color{gray},
  frame=single, breaklines=true, showstringspaces=false,
  language=Python
}

% ------------------------------------------------------------------
\title{\textbf{Algoritma Shor dan Ancamannya terhadap Kriptosistem RSA:\\
Analisis Teoretis dan Simulasi dengan Qiskit}}
\author{Calvin Sabastian Tanzil (10223025)\\
Muhammad Fatih Irkham Mauludi (13524004)\\[2pt]
\small Institut Teknologi Bandung\\
\small FI4111 Sistem Kuantum Terbuka dan Komputasi Kuantum}
\date{\today}

\begin{document}
\maketitle

\begin{abstract}
Keamanan kriptosistem RSA bertumpu pada asumsi bahwa memfaktorkan bilangan
bulat besar $N = pq$ tidak dapat dilakukan secara efisien oleh komputer klasik.
Algoritma Shor menunjukkan bahwa komputer kuantum dapat menyelesaikan masalah
tersebut dalam waktu polinomial terhadap jumlah bit $N$. Makalah ini membahas
landasan matematis RSA, prinsip komputasi kuantum yang relevan, serta
mekanisme dan arsitektur rangkaian algoritma Shor, kemudian
mengimplementasikannya dalam simulator Qiskit untuk kasus RSA
\emph{toy-sized}. \todo{lengkapi ringkasan hasil simulasi dan kesimpulan utama setelah Bab 4 selesai.}

\medskip
\noindent\textbf{Kata kunci:} RSA, algoritma Shor, estimasi fase kuantum,
transformasi Fourier kuantum, Qiskit.
\end{abstract}

% ==================================================================
\section{Pendahuluan}
% ==================================================================

\subsection{Latar Belakang}
Kriptografi kunci publik merupakan fondasi keamanan komunikasi digital modern,
mulai dari protokol TLS pada peramban web hingga tanda tangan digital pada
dokumen elektronik. Salah satu skema kunci publik yang paling luas digunakan
adalah RSA, yang diperkenalkan oleh Rivest, Shamir, dan Adleman pada tahun
1977 \cite{rsa1978}. Keamanan RSA bergantung pada sulitnya memfaktorkan
modulus $N$ yang merupakan hasil kali dua bilangan prima besar. Algoritma
klasik terbaik yang dikenal, yaitu \emph{General Number Field Sieve} (GNFS),
memiliki kompleksitas subeksponensial
\begin{equation}
  T_{\text{GNFS}}(N) = \exp\!\left[\left(\tfrac{64}{9}\right)^{1/3}
  (\ln N)^{1/3} (\ln\ln N)^{2/3}\,(1+o(1))\right],
\end{equation}
sehingga untuk ukuran kunci 2048 bit, faktorisasi secara klasik dianggap tidak
praktis.

Pada tahun 1994, Peter Shor menunjukkan bahwa sebuah komputer kuantum dapat
memfaktorkan $N$ dalam waktu polinomial, kira-kira $O\big((\log N)^3\big)$
operasi gate \cite{shor1997}. Hasil ini mengubah pemahaman tentang
batas keamanan kriptografi kunci publik dan menjadi salah satu motivasi utama
pengembangan komputer kuantum sekaligus kriptografi pascakuantum
(\emph{post-quantum cryptography}) \cite{nist_pqc}. Walaupun komputer kuantum
saat ini belum mampu menjalankan algoritma Shor pada ukuran kunci yang
relevan, memahami cara kerja algoritma ini penting untuk menilai urgensi
migrasi ke skema kriptografi yang tahan kuantum.

\subsection{Rumusan Masalah}
\begin{enumerate}[label=\arabic*.]
  \item Bagaimana landasan matematis kriptosistem RSA dan mengapa keamanannya
        bergantung pada masalah faktorisasi?
  \item Bagaimana mekanisme dan arsitektur rangkaian kuantum algoritma Shor,
        termasuk konstruksi \emph{oracle}-nya?
  \item Bagaimana hasil simulasi algoritma Shor menggunakan Qiskit untuk kasus
        RSA berukuran kecil, dan apa yang dapat disimpulkan dari hasil tersebut?
\end{enumerate}

\subsection{Tujuan}
\begin{enumerate}[label=\arabic*.]
  \item Menjelaskan dasar teori bilangan yang melandasi RSA.
  \item Menguraikan mekanisme algoritma Shor beserta arsitektur rangkaiannya.
  \item Mengimplementasikan dan menganalisis simulasi algoritma Shor
        menggunakan Qiskit.
\end{enumerate}

\subsection{Batasan Masalah}
\todo{sesuaikan dengan scope yang kamu pilih.} Contoh batasan:
\begin{enumerate}[label=\arabic*.]
  \item Kasus yang disimulasikan adalah RSA \emph{toy-sized}
        (misalnya $N = 15$), bukan ukuran kunci yang digunakan dalam praktik.
  \item Simulasi dijalankan pada simulator klasik (\texttt{qiskit-aer}), bukan
        pada perangkat keras kuantum.
  \item \todo{apakah simulasi dengan model derau (\emph{noise}) termasuk dalam cakupan?}
  \item \todo{apakah konstruksi oracle bersifat khusus untuk $N$ tertentu atau umum?}
\end{enumerate}

\subsection{Sistematika Penulisan}
Bab~\ref{sec:teori} memaparkan landasan teori RSA dan komputasi kuantum.
Bab~\ref{sec:mekanisme} menguraikan mekanisme dan arsitektur algoritma Shor.
Bab~\ref{sec:implementasi} menyajikan implementasi Qiskit beserta hasil dan
pembahasannya. Bab~\ref{sec:kesimpulan} berisi kesimpulan dan saran.

% ==================================================================
\section{Landasan Teori}\label{sec:teori}
% ==================================================================

\subsection{Teori Bilangan untuk RSA}

\begin{definisi}[Kongruensi]
Untuk bilangan bulat $a$, $b$, dan $n > 0$, ditulis $a \equiv b \pmod n$
jika $n \mid (a - b)$.
\end{definisi}

\begin{definisi}[\emph{Euler Totient Function}]
$\varphi(n)$ adalah banyaknya bilangan bulat $1 \le k \le n$ yang memenuhi
$\gcd(k, n) = 1$. Untuk $n = pq$ dengan $p \ne q$ prima,
\begin{equation}
  \varphi(pq) = (p-1)(q-1).
\end{equation}
\end{definisi}

\begin{teorema}[Euler]
Jika $\gcd(a, n) = 1$, maka $a^{\varphi(n)} \equiv 1 \pmod n$.
\end{teorema}

\begin{definisi}[Invers Modular]
Jika $\gcd(e, m) = 1$, terdapat $d$ yang unik modulo $m$ sehingga
$ed \equiv 1 \pmod m$. Nilai $d$ dapat dihitung secara efisien dengan
algoritma Euclid yang diperluas.
\end{definisi}

\subsection{Kriptosistem RSA}

\paragraph{Pembangkitan kunci.}
\begin{enumerate}[label=\arabic*.]
  \item Pilih dua bilangan prima berbeda $p$ dan $q$, lalu hitung $N = pq$.
  \item Hitung $\varphi(N) = (p-1)(q-1)$.
  \item Pilih eksponen publik $e$ dengan $1 < e < \varphi(N)$ dan
        $\gcd(e, \varphi(N)) = 1$.
  \item Hitung eksponen privat $d \equiv e^{-1} \pmod{\varphi(N)}$.
\end{enumerate}
Kunci publik adalah $(N, e)$, sedangkan kunci privat adalah $d$
(beserta $p$ dan $q$ yang dirahasiakan).

\paragraph{Enkripsi dan dekripsi.}
Untuk pesan $m$ dengan $0 \le m < N$,
\begin{equation}
  c = m^e \bmod N, \qquad m = c^d \bmod N.
\end{equation}

\paragraph{Kebenaran.}
Karena $ed = 1 + k\varphi(N)$ untuk suatu bilangan bulat $k$, untuk
$\gcd(m, N) = 1$ berlaku
\begin{equation}
  c^d \equiv m^{ed} = m \cdot \big(m^{\varphi(N)}\big)^k \equiv m \pmod N
\end{equation}
berdasarkan teorema Euler. (Kasus $\gcd(m, N) \ne 1$ dapat dibuktikan dengan
\emph{Chinese Remainder Theorem}.)

\paragraph{Hubungan dengan faktorisasi.}
Siapa pun yang mengetahui $p$ dan $q$ dapat menghitung $\varphi(N)$ lalu
memperoleh $d$ dengan algoritma Euclid yang diperluas. Dengan demikian,
kemampuan memfaktorkan $N$ secara efisien cukup untuk mematahkan RSA. Inilah
alasan mengapa algoritma faktorisasi yang efisien, baik klasik maupun kuantum,
merupakan ancaman langsung terhadap RSA.

\subsection{Dasar Komputasi Kuantum}

\subsubsection{Qubit dan Register}
Sebuah qubit adalah sistem kuantum dua tingkat dengan keadaan umum
\begin{equation}
  \ket{\psi} = \alpha\ket{0} + \beta\ket{1}, \qquad |\alpha|^2 + |\beta|^2 = 1.
\end{equation}
Register $n$ qubit hidup dalam ruang Hilbert berdimensi $2^n$ dengan basis
komputasi $\{\ket{x} : x = 0, 1, \dots, 2^n - 1\}$.

\subsubsection{Gerbang Kuantum Dasar}
Gerbang Hadamard menghasilkan superposisi seragam:
\begin{equation}
  H = \frac{1}{\sqrt 2}\begin{pmatrix} 1 & 1 \\ 1 & -1 \end{pmatrix},
  \qquad
  H^{\otimes n}\ket{0}^{\otimes n} = \frac{1}{\sqrt{2^n}} \sum_{x=0}^{2^n-1} \ket{x}.
\end{equation}
Gerbang fase $P(\phi) = \mathrm{diag}(1, e^{i\phi})$ beserta versi
terkendalinya merupakan komponen utama transformasi Fourier kuantum.

\subsubsection{Transformasi Fourier Kuantum (QFT)}
QFT pada $n$ qubit didefinisikan sebagai
\begin{equation}
  \mathrm{QFT}\ket{x} = \frac{1}{\sqrt{2^n}} \sum_{k=0}^{2^n-1}
  e^{2\pi i\, x k / 2^n} \ket{k}.
\end{equation}
QFT dapat diimplementasikan dengan $O(n^2)$ gerbang (Hadamard dan fase
terkendali), jauh lebih efisien dibanding FFT klasik pada vektor berdimensi
$2^n$ yang membutuhkan $O(n 2^n)$ operasi \cite{nielsen2010}.

\subsubsection{Estimasi Fase Kuantum (QPE)}
Misalkan $U$ adalah operator uniter dengan vektor eigen $\ket{u}$ dan nilai
eigen $e^{2\pi i \theta}$, $\theta \in [0, 1)$. QPE menggunakan register
penghitung berisi $t$ qubit untuk mengestimasi $\theta$:
\begin{enumerate}[label=\arabic*.]
  \item Siapkan $\ket{0}^{\otimes t}\ket{u}$, lalu terapkan $H^{\otimes t}$
        pada register penghitung.
  \item Terapkan $U^{2^j}$ terkendali oleh qubit penghitung ke-$j$, sehingga
        diperoleh
        \begin{equation}
          \frac{1}{\sqrt{2^t}} \sum_{k=0}^{2^t-1} e^{2\pi i \theta k}\ket{k}\ket{u}.
        \end{equation}
  \item Terapkan $\mathrm{QFT}^\dagger$ pada register penghitung, lalu ukur.
\end{enumerate}
Hasil pengukuran $\tilde y$ memberikan estimasi $\theta \approx \tilde y / 2^t$.
Jika $\theta\, 2^t$ adalah bilangan bulat, estimasinya eksak; jika tidak,
peluang memperoleh aproksimasi terbaik sedikitnya $4/\pi^2 \approx 0{,}405$
\cite{nielsen2010}.

\subsection{Algoritma Shor secara Garis Besar}
Algoritma Shor memadukan bagian klasik dan bagian kuantum. Bagian kuantumnya
adalah penerapan QPE untuk mencari \emph{orde} suatu bilangan modulo $N$,
sedangkan bagian klasiknya mengubah orde tersebut menjadi faktor $N$ melalui
perhitungan FPB (faktor persekutuan terbesar). Rincian mekanisme dan
arsitekturnya dibahas pada Bab~\ref{sec:mekanisme}.

\subsection{Qiskit}
Qiskit adalah \emph{software development kit} sumber terbuka untuk komputasi
kuantum yang dikembangkan IBM \cite{qiskit2024}. Rangkaian dibangun dengan
kelas \texttt{QuantumCircuit}, ditranspilasi ke himpunan gerbang dasar dengan
\texttt{transpile}, lalu dijalankan pada simulator \texttt{AerSimulator}
dari paket \texttt{qiskit-aer} atau pada perangkat keras IBM Quantum.
\todo{sebutkan versi Qiskit dan qiskit-aer yang dipakai.}

% ==================================================================
\section{Mekanisme dan Arsitektur Algoritma}\label{sec:mekanisme}
% ==================================================================

\todo{Bagian ini diisi sendiri.} Saran struktur:

\subsection{Reduksi Faktorisasi ke Pencarian Orde}
\todo{definisi orde $r$; syarat $r$ genap dan $a^{r/2} \not\equiv -1 \pmod N$;
bagaimana faktor diperoleh dari FPB; peluang keberhasilan untuk $a$ acak.}

\subsection{Operator Uniter dan Struktur Nilai Eigennya}
\todo{definisi $U_a$; vektor eigen $\ket{u_s}$ dan nilai eigennya; mengapa
$\ket{1}$ dapat dipakai sebagai keadaan awal register kerja.}

\subsection{Konstruksi Oracle Eksponensiasi Modular}
\todo{konstruksi oracle yang kamu pilih; diagram rangkaian; jumlah qubit dan
kedalaman rangkaian.}

\subsection{Arsitektur Rangkaian Lengkap}
\todo{register penghitung dan register kerja, urutan gerbang, pengukuran.
Sisipkan gambar rangkaian dari Qiskit:}
% \begin{figure}[h]
%   \centering
%   \includegraphics[width=\textwidth]{../figures/rangkaian.pdf}
%   \caption{Rangkaian lengkap ...}
%   \label{fig:rangkaian}
% \end{figure}

\subsection{Pascaproses Klasik: Pecahan Berlanjut}
\todo{dari hasil ukur ke estimasi fase, lalu ke orde $r$.}

% ==================================================================
\section{Implementasi Qiskit, Simulasi, dan Pembahasan}\label{sec:implementasi}
% ==================================================================

\subsection{Lingkungan Simulasi}
\todo{versi Python, Qiskit, qiskit-aer; jumlah \emph{shots}; \emph{seed}.}

\subsection{Parameter RSA yang Digunakan}
\todo{isi tabel sesuai parameter yang kamu pakai.}
\begin{table}[h]
  \centering
  \caption{Parameter RSA \emph{toy-sized}.}
  \label{tab:rsa}
  \begin{tabular}{lc}
    \toprule
    Besaran & Nilai \\
    \midrule
    $p, q$        & \todo{} \\
    $N$           & \todo{} \\
    $\varphi(N)$  & \todo{} \\
    $e$           & \todo{} \\
    $d$           & \todo{} \\
    \bottomrule
  \end{tabular}
\end{table}

\subsection{Implementasi Rangkaian}
\todo{cuplikan kode penting. Contoh memasukkan kode dari berkas:}
% \lstinputlisting[firstline=1,lastline=40,caption={...}]{../src/nama_berkas.py}

\subsection{Hasil Simulasi}
\todo{histogram hasil pengukuran, tabel fase terukur dan orde kandidat.}
% \begin{figure}[h]
%   \centering
%   \includegraphics[width=0.8\textwidth]{../figures/histogram.pdf}
%   \caption{Distribusi hasil pengukuran register penghitung.}
%   \label{fig:hist}
% \end{figure}

\subsection{Pembahasan}
\todo{interpretasi puncak histogram; keberhasilan per pilihan basis;
sumber kegagalan; kedalaman rangkaian setelah transpilasi; perbandingan
dengan derau (jika dalam cakupan).}

\subsection{Implikasi terhadap Keamanan RSA}
Estimasi sumber daya terbaru menunjukkan bahwa mematahkan RSA-2048 dengan
algoritma Shor membutuhkan komputer kuantum dengan koreksi galat berskala
sangat besar, jauh melampaui kemampuan perangkat saat ini
\cite{gidney2025}. \todo{verifikasi angka terbaru dari literatur dan
kaitkan dengan hasil simulasimu.} Meskipun demikian, ancaman
\emph{harvest now, decrypt later} membuat migrasi ke kriptografi
pascakuantum, seperti standar ML-KEM dan ML-DSA dari NIST \cite{nist_pqc},
menjadi agenda yang mendesak.

% ==================================================================
\section{Kesimpulan}\label{sec:kesimpulan}
% ==================================================================

\subsection{Kesimpulan}
\begin{enumerate}[label=\arabic*.]
  \item Keamanan RSA bergantung pada sulitnya faktorisasi $N = pq$; dengan
        mengetahui $p$ dan $q$, kunci privat $d$ dapat dihitung secara
        efisien.
  \item \todo{kesimpulan tentang mekanisme dan arsitektur algoritma Shor.}
  \item \todo{kesimpulan dari hasil simulasi Qiskit.}
\end{enumerate}

\subsection{Saran}
\todo{misalnya: simulasi untuk $N$ yang lebih besar, eksekusi pada perangkat
keras IBM Quantum, kajian mitigasi galat, atau perbandingan dengan
kriptografi pascakuantum.}

% ==================================================================
\begin{thebibliography}{9}

\bibitem{rsa1978}
R.~L. Rivest, A.~Shamir, dan L.~Adleman,
``A method for obtaining digital signatures and public-key cryptosystems,''
\emph{Communications of the ACM}, vol.~21, no.~2, pp.~120--126, 1978.

\bibitem{shor1997}
P.~W. Shor,
``Polynomial-time algorithms for prime factorization and discrete logarithms
on a quantum computer,''
\emph{SIAM Journal on Computing}, vol.~26, no.~5, pp.~1484--1509, 1997.

\bibitem{nielsen2010}
M.~A. Nielsen dan I.~L. Chuang,
\emph{Quantum Computation and Quantum Information}, edisi ke-10.
Cambridge University Press, 2010.

\bibitem{qiskit2024}
A.~Javadi-Abhari \emph{et al.},
``Quantum computing with Qiskit,'' arXiv:2405.08810, 2024.

\bibitem{gidney2025}
C.~Gidney,
``How to factor 2048 bit RSA integers with less than a million noisy qubits,''
arXiv:2505.15917, 2025.

\bibitem{nist_pqc}
National Institute of Standards and Technology,
\emph{FIPS 203: Module-Lattice-Based Key-Encapsulation Mechanism Standard},
2024.

\end{thebibliography}

\end{document}
